\documentclass[11pt]{article}

\usepackage[margin=1in]{geometry}
\usepackage{amsmath,amssymb}
\usepackage{array}
\usepackage{booktabs}
\usepackage{hyperref}
\usepackage{xcolor}
\usepackage{enumitem}
\usepackage{listings}

\hypersetup{
  colorlinks=true,
  linkcolor=blue!60!black,
  urlcolor=blue!60!black,
  citecolor=blue!60!black
}

\newcommand{\code}[1]{\texttt{\detokenize{#1}}}
\newcommand{\R}{\mathbb{R}}
\newcommand{\E}{\mathcal{E}}
\newcommand{\K}{\mathcal{K}}
\newcommand{\G}{\mathcal{G}}
\newcommand{\V}{\mathcal{V}}
\newcommand{\A}{\mathcal{A}}

\lstset{
  basicstyle=\ttfamily\small,
  columns=fullflexible,
  keepspaces=true,
  frame=single,
  breaklines=true,
  backgroundcolor=\color{black!3}
}

\title{Upwind Graph Ordering Algorithm}
\author{\code{hdgfem.linalg.ordering}}
\date{}

\begin{document}
\maketitle

\begin{abstract}
This note documents the ordering algorithm implemented in
\code{linalg/ordering.py}.  The purpose of the module is to build an
experimental permutation for HDG trace unknowns in advection-dominated systems.
The ordering is algebraic with respect to the sparse trace system, but it is
constructed from mesh-edge connectivity and the sign of the normal advective
flux on element faces.  The main pipeline is:
\begin{enumerate}[nosep]
  \item build a directed graph between mesh-edge blocks using local
        inflow-to-outflow dependencies;
  \item build forward CSR and try deterministic Kahn ordering;
  \item for an acyclic graph, use that order directly;
  \item otherwise trim acyclic prefixes and suffixes, run Kosaraju only on the
        compact residual, and topologically order the recombined components;
  \item lift the ordered edge blocks to a trace degree-of-freedom permutation.
\end{enumerate}
The module also contains sparsity-pattern plotting utilities for comparing the
matrix before and after the upwind SCC permutation.
\end{abstract}

\tableofcontents

\section{Scope and Design Intent}

The HDG global unknown is the trace field on mesh edges.  For a uniform trace
polynomial degree, each global mesh edge carries the same number of trace
degrees of freedom, denoted here by
\[
  n_{\partial} = \code{edg_dof}.
\]
The graph-ordering module does not alter element assembly, mesh connectivity,
or numerical fluxes.  It only computes a permutation of trace unknown indices.
The convention matches \code{hdgfem.linalg.system.solve_global_system}:
\[
  A_{\pi} = A[\pi, \pi], \qquad b_{\pi} = b[\pi],
\]
and after solving the permuted system the solution is unpermuted by
\[
  x[\pi] = x_{\pi}.
\]

The ordering is designed to expose approximate upstream-to-downstream
structure in advection-dominated HDG trace matrices.  It should be viewed as a
preconditioner-ordering experiment, not as a guaranteed fill-reducing ordering.
In particular, SuperLU's internal orderings such as \code{COLAMD} may still be
better at high polynomial order because they optimize symbolic fill rather than
flow direction.

\section{Inputs and Core Data Structures}

\subsection{Mesh edge connectivity}

The primary mesh input is the element-local to global-edge map
\[
  \code{loc2glob_edge}[K,f] \in \{0,\ldots,N_e-1\},
\]
where \(K\) is an element index, \(f\in\{0,1,2\}\) is a local triangle face, and
\(N_e\) is the number of global mesh edges.  In the code this is passed to the
Numba kernel as \code{loc2glob_edge}.  For boundary-eliminated trace systems,
only non-boundary edge blocks remain active.  The helper \code{_as_active_edges}
builds three arrays:
\begin{itemize}[nosep]
  \item \code{active_edge_ids}: sorted global edge ids participating in the
        reduced system;
  \item \code{active_mask}: Boolean mask over all mesh edges;
  \item \code{old_to_active}: map from global edge id to active graph node id,
        with \(-1\) for inactive edges.
\end{itemize}

\subsection{Face-normal advective flux}

The second geometric/numerical input is
\[
  \code{beta_dot_normal}[K,f,q]
  \approx \beta_h(x_{K,f,q})\cdot n_{K,f},
\]
where \(q\) indexes face quadrature points.  The graph construction reduces
this to one mean face flux per element-local face:
\[
  \bar F_{K,f}
  =
  \frac{1}{N_q}\sum_{q=1}^{N_q}
  \beta_h(x_{K,f,q})\cdot n_{K,f}.
\]
A face is classified as inflow when
\[
  \bar F_{K,f} < -\varepsilon,
\]
and as outflow when
\[
  \bar F_{K,f} > \varepsilon,
\]
where \(\varepsilon=\code{flux_tolerance}\).  Faces in the neutral band
\([-\varepsilon,\varepsilon]\) are ignored.

\subsection{Return containers}

The algorithm returns a \code{GraphOrderingResult}:
\[
\begin{array}{ll}
\code{edge_order} & \text{active global edge ids in graph order},\\
\code{dof_permutation} & \text{trace dof permutation used by the sparse solver},\\
\code{component_id} & \text{SCC id for each active graph node},\\
\code{component_order} & \text{topological order of SCC ids},\\
\code{diagnostics} & \text{graph sizes and timing information}.
\end{array}
\]
The diagnostics contain the number of graph nodes, directed graph edges, SCCs,
the largest SCC size, the number of cyclic SCCs, the number of nodes in cyclic
SCCs, and a timing breakdown.

\section{Step 1: Build the Directed Upwind Edge Graph}

\subsection{Graph nodes}

The directed graph nodes are active mesh-edge blocks, not individual trace
degrees of freedom.  If boundary dofs have not been eliminated, every mesh edge
is active.  If boundary dofs have been eliminated, only non-boundary edges are
active.  Therefore the graph vertex set is
\[
  \V = \{0,\ldots,N_a-1\},
\]
where \(N_a=\code{active_edge_ids.size}\).

\subsection{Element-local dependencies}

For each element \(K\), the kernel inspects all local faces.  Every active
inflow face is connected to every active outflow face in the same element:
\[
  e_{K,f_{\rm in}} \longrightarrow e_{K,f_{\rm out}}.
\]
In code, this is done by \code{_build_upwind_edge_pairs_kernel}.  In words:

\begin{lstlisting}
for each element K:
    compute mean face flux F[f] for f=0,1,2
    for each source face f:
        keep it only if active and F[f] < -flux_tolerance
        for each target face g != f:
            keep it only if active and F[g] > flux_tolerance
            append edge active_id(f) -> active_id(g)
\end{lstlisting}

The direction ``inflow to outflow'' is the local transport dependency direction:
information entering an element through an inflow face contributes to unknowns
on outflow faces.  The graph can contain duplicate directed edges.  Duplicates
are accepted because they do not affect reachability or SCC membership, and
removing them would add a sorting/unique phase.

\subsection{Storage}

The kernel emits two arrays
\[
  \code{sources},\quad \code{targets}
\]
with equal length.  Directed graph edge \(m\) is
\[
  \code{sources}[m]\rightarrow \code{targets}[m].
\]
The pair builder is a deterministic two-pass parallel algorithm.  A first
\code{prange} pass computes the three mean face fluxes and the exact pair count
for every element.  A serial prefix sum assigns each element a disjoint output
slice.  A second \code{prange} pass fills that slice in local source-face then
target-face order.  Thus the maximum remains \(6N_K\), but allocation uses the
exact pair count and the result is invariant under thread scheduling.

\section{Step 2: Forward CSR and Acyclic Detection}

The helper \code{_build_csr_kernel} converts edge arrays into compressed sparse
row adjacency:
\[
  \code{indptr}\in\mathbb{N}^{N_a+1}, \qquad
  \code{indices}\in\mathbb{N}^{M_g}.
\]
Only this forward CSR is built initially.  The implementation counts outgoing
edges, prefix-sums those counts, and fills adjacency entries without removing
duplicates.

The cached \code{_kahn_frontier_kernel} computes indegrees and processes all
currently available zero-indegree nodes as one deterministic frontier.  Each
new frontier is sorted by node id.  When the frontier-node plus adjacency-visit
work is at least 32,768, copying the frontier adjacency uses \code{prange};
smaller frontiers stay serial to avoid scheduler overhead.  Indegree updates
are always applied in deterministic packed-buffer order, so atomics and
thread-dependent tie breaking are unnecessary.

If Kahn emits all \(N_a\) nodes, the graph is a DAG.  This
\code{acyclic-fast} path uses the Kahn order directly, assigns one singleton
component to each node, and bypasses reverse CSR, Kosaraju, component-edge
construction, and a second topological sort.

\section{Step 3: Cyclic Residual Processing}

If the first Kahn pass stops early, its emitted nodes are the maximal acyclic
prefix.  The algorithm then builds reverse CSR and runs the same frontier
kernel on the still-active nodes.  Zero-indegree nodes of the reverse graph are
zero-outdegree nodes of the original graph; recursively removing them peels the
acyclic suffix.  Isolated and disconnected acyclic nodes are peeled as well.

The remaining mask defines the compact residual.  Residual nodes are mapped to
consecutive ids in ascending original-node order, and only edges whose two ends
remain active are copied.  If the full graph has \(N_a,M_g\) nodes and edges,
write the residual sizes as \(N_r,M_r\).

The function \code{_kosaraju_scc_kernel} runs only on this residual.  It uses
explicit stacks rather than Python recursion.  Its first DFS pass records
forward-graph finishing order and its second pass visits reverse CSR in reverse
finishing order.  It returns
\[
  \code{residual_component_id}[v] = c,
  \qquad
  \code{residual_component_sizes}[c] = |C_c|.
\]
Kosaraju therefore costs \(O(N_r+M_r)\), not \(O(N_a+M_g)\), when trimming is
effective.  A singleton self-loop remains in the residual and is reported as a
cyclic component even though its component size is one.

\section{Step 4: Recombine and Order Components}

The function \code{_recombine_components_kernel} maps residual SCCs back to
original nodes and assigns every peeled node its own singleton component.
Scanning original nodes in ascending order assigns component ids by their
smallest original node, independent of worker count and residual SCC labels.

The function \code{_component_edges_kernel} maps each original edge
\(u\rightarrow v\) to \(C(u)\rightarrow C(v)\), discarding component
self-edges.  The result is a DAG containing both residual SCCs and peeled
singletons.  A final \code{_kahn_frontier_kernel} call orders all components.
The emitted component count must equal the number of components; otherwise the
SCC partition or recombination is invalid.

\section{Step 5: Lift Component Order to Edge-Block Order}

The function \code{_node_order_from_components_kernel} produces an ordered list
of graph nodes.  It reserves contiguous ranges for components according to the
topological component order, then inserts original node ids into the range of
their component.  Thus:
\begin{itemize}[nosep]
  \item components are topologically ordered;
  \item nodes inside one SCC keep original active-node order.
\end{itemize}
The active-node order is then converted back to global mesh-edge ids:
\[
  \code{edge_order}
  =
  \code{active_edge_ids}[\code{node_order}].
\]

\section{Step 6: Lift Edge Blocks to Trace DOFs}

The sparse trace matrix is indexed by individual trace degrees of freedom.  If
edge block \(e\) has \(n_{\partial}\) local trace dofs, its trace dof range is
\[
  e n_{\partial},\ e n_{\partial}+1,\ \ldots,\ e n_{\partial}+n_{\partial}-1.
\]
The helper \code{_trace_dof_permutation_kernel} receives ordered active-node
positions, not global edge ids.  For each ordered active node position \(a\),
it appends the contiguous local block
\[
  a n_{\partial},\ a n_{\partial}+1,\ldots,a n_{\partial}+n_{\partial}-1.
\]
This is exactly the correct indexing for reduced systems where active edge
blocks have been compacted into consecutive reduced-system coordinates.

For a full, non-reduced trace system, active nodes are all edges in natural
global edge order.  Then the active-node position and global edge id coincide.
For a boundary-eliminated system, \code{active_edges} must be passed so the
active edge positions match the reduced-system matrix.

\section{Boundary Elimination and Active Edges}

The ordering routine is used in two common modes:

\paragraph{Full penalty system.}
Boundary trace dofs remain in the system.  The active edge set is all mesh
edges:
\[
  \code{active_edges} = \code{None}.
\]

\paragraph{Boundary-eliminated system.}
Prescribed boundary trace dofs are removed before the solve.  The reduced
matrix contains only non-boundary edge blocks.  In this case the caller builds
\[
  \code{active_edges} = \{e:\ e \notin \partial\Omega\}.
\]
The graph is built using global edge ids, but all graph nodes are compacted
through \code{old_to_active}.  The resulting dof permutation is therefore in
reduced-system coordinates and can be applied directly to the reduced matrix.

\section{Complexity}

Let \(N_K\) be the number of elements, \(N_a\) the number of active edge
blocks, and \(M_g\) the number of directed graph edges after local inflow to
outflow construction.  For triangular elements,
\[
  M_g \le 6N_K.
\]
The common acyclic path has costs
\[
\begin{array}{ll}
\text{graph pair construction} & O(N_K N_q),\\
\text{forward CSR and Kahn order} & O(N_a + M_g),\\
\text{level diagnostics} & O(N_a + M_g),\\
\text{dof lift} & O(N_a n_{\partial}).
\end{array}
\]
It stores only forward CSR.  For a cyclic graph, reverse CSR and trimming add
\(O(N_a+M_g)\); compact residual CSR and Kosaraju add \(O(N_r+M_r)\); and
recombined condensation construction and ordering add
\(O(M_g+N_{\rm comp}+M_{\rm comp})\).  Duplicate edges are retained and are
included in these edge counts.

The graph-ordering cost is usually small compared with HDG matrix assembly and
incomplete factorization.  It is also almost independent of the polynomial
degree \(p\), except for the final dof lift through \(n_{\partial}=p+1\).

All Numba kernels use \code{cache=True}.  Common kernels are compiled or loaded
by both paths, while compaction, Kosaraju, recombination, and component-edge
kernels are reached lazily only for cyclic inputs.  Pair and large-frontier
loops use the active Numba thread pool; their fixed output slices and sorted
frontiers keep results independent of its size.

\section{Diagnostics}

\code{GraphOrderingDiagnostics} reports:
\[
\begin{array}{ll}
\code{num_nodes} & N_a,\\
\code{num_directed_edges} & M_g,\\
\code{num_components} & N_{\rm comp},\\
\code{largest_component_size} & \max_c |C_c|,\\
\code{cyclic_components} & \text{multi-node SCCs plus singleton self-loops},\\
\code{cyclic_nodes} & \text{nodes in those cyclic components},\\
\code{algorithm_path} & \code{acyclic-fast}\text{ or }\code{cyclic-residual},\\
\code{peeled_nodes} & \text{trimmed nodes on the cyclic path},\\
\code{residual_nodes} & N_r\text{ on the cyclic path}.
\end{array}
\]
For a nearly acyclic upwind graph, one typically sees
\[
  \code{largest_component_size}=1,\qquad \code{cyclic_nodes}=0.
\]
For recirculating flows, SCCs larger than one indicate trace-edge blocks whose
flow dependencies are mutually coupled.

The timing object \code{GraphOrderingTimings} retains graph-pair, CSR, SCC,
DAG, topological-order, dof-permutation, and total times.  The topological field
now measures only the successful Kahn sort.  Separate \code{level_diagnostics}
and \code{node_order} fields measure the compiled level-width and node-lifting
phases.

\section{Sparse Pattern Plotting}

The module also provides two plotting helpers:
\[
\begin{array}{ll}
\code{save_sparse_pattern_plot} & \text{plot one sparse matrix pattern},\\
\code{save_upwind_reordered_matrix_patterns}
  & \text{plot before and after } A[\pi,\pi].
\end{array}
\]
The marker area is chosen by \code{sparse_pattern_marker_area}.  The intended
law is inverse in the number of nonzeros:
\[
  s
  =
  s_{\rm ref}\frac{N_{\rm ref}}{\operatorname{nnz}(A)}.
\]
The current implementation clips this value and applies an additional practical
floor, so very large matrices do not disappear completely in PNG output.  Very
large matrices are also stride-sampled to \code{max_plot_points}; the marker
area is still computed from the original \(\operatorname{nnz}(A)\), not from
the sampled count.

\section{Interaction with SuperLU ILU}

The upwind SCC permutation is not a symbolic fill-reducing ordering.  It is a
physics-motivated block ordering.  If the sparse solver is SuperLU's
\code{spilu}, then its \code{permc_spec} option matters:
\begin{itemize}
  \item \code{NATURAL} forces SuperLU to respect the supplied matrix ordering.
        This can expose the benefit of a good upwind order, but it can also be
        fragile at high polynomial degree.
  \item \code{COLAMD} lets SuperLU apply a fill-reducing column ordering.  This
        may override much of the upwind-ordering effect, but it can be far more
        robust when local trace blocks become large.
\end{itemize}
Thus the SCC ordering should be benchmarked together with the ILU column
permutation strategy.  A small edge-block graph can look excellent while the
lifted high-order dof matrix still has poor symbolic fill under \code{NATURAL}.

\section{High-Level Pseudocode}

\begin{lstlisting}
sources, targets = parallel_count_prefix_fill(...)
forward_csr = build_csr(sources, targets)
prefix_or_order = kahn_frontiers(forward_csr)

if size(prefix_or_order) == num_active_nodes:
    algorithm_path = "acyclic-fast"
    component_id = arange(num_active_nodes)
    component_order = prefix_or_order
else:
    algorithm_path = "cyclic-residual"
    reverse_csr = build_csr(targets, sources)
    suffix = kahn_frontiers(reverse_csr, excluding=prefix_or_order)
    residual = compact(excluding=prefix_or_order + suffix)
    residual_components = kosaraju(residual.forward, residual.reverse)
    component_id = recombine(residual_components,
                             peeled_singletons=prefix_or_order + suffix)
    component_dag = build_component_graph(component_id, sources, targets)
    component_order = kahn_frontiers(component_dag)

node_order = nodes_by_component(component_id, component_order)
level_widths = compiled_level_diagnostics(component_dag, component_order)
dof_permutation = lift_edge_order_to_trace_dofs(node_order, edg_dof)
return GraphOrderingResult(..., diagnostics.algorithm_path=algorithm_path)
\end{lstlisting}

\section{Function Map}

\begin{center}
\small
\begin{tabular}{p{0.42\linewidth}p{0.48\linewidth}}
\toprule
Function & Role \\
\midrule
\code{_build_upwind_edge_pairs_kernel} & Parallel deterministic pair count and fill. \\
\code{_build_csr_kernel} & Convert directed edge arrays to CSR adjacency. \\
\code{_kahn_frontier_kernel} & Detect a DAG, trim frontiers, or order components. \\
\code{_compact_residual_graph_kernel} & Compact cyclic residual nodes and edges. \\
\code{_kosaraju_scc_kernel} & Compute residual SCCs with iterative Kosaraju DFS. \\
\code{_recombine_components_kernel} & Merge residual SCCs and peeled singletons. \\
\code{_component_edges_kernel} & Build condensation graph edges between SCCs. \\
\code{_level_widths_kernel} & Compute DAG levels and width statistics. \\
\code{_node_order_from_components_kernel} & Convert component order to node order. \\
\code{_trace_dof_permutation_kernel} & Lift ordered edge blocks to trace dofs. \\
\code{strongly_connected_component_order} & Public adaptive graph-only routine. \\
\code{upwind_scc_trace_ordering} & Public HDG trace-edge ordering routine. \\
\code{sparse_pattern_marker_area} & Compute nnz-dependent marker area for plots. \\
\code{save_sparse_pattern_plot} & Save one matrix sparsity pattern figure. \\
\code{save_upwind_reordered_matrix_patterns} & Save before/after upwind permutation figures. \\
\bottomrule
\end{tabular}
\end{center}

\section{Practical Usage}

The maintained advection-reaction case runner can generate matrix patterns
without entering the global solve:

\begin{lstlisting}
.venv/bin/python scripts/advection_reaction/run_cases.py \
  -p 7 --lc 0.03 --boundary-mode eliminate \
  --trace-ordering upwind-scc --matrix-pattern-only \
  --matrix-pattern-max-points 2000000
\end{lstlisting}

This writes figures such as:
\begin{lstlisting}
matrix_patterns/*_before_upwind.png
matrix_patterns/*_after_upwind.png
\end{lstlisting}

\section{Summary}

The ordering in \code{hdgfem.linalg.ordering} is a block-level flow ordering:
construct local inflow-to-outflow dependencies in parallel, use forward-only
Kahn ordering for DAGs, and restrict SCC work to a trimmed cyclic residual
before recombining and ordering components.  The result is lifted to trace
dofs.  The graph construction is inexpensive and mostly independent of
polynomial order.  Its effect on ILU depends on how the block ordering interacts
with the symbolic fill behavior of the sparse factorization.

\end{document}
